% !TeX program = xelatex
% ====================================================================
% yuhang_yang_cv.tex -- Chinese academic CV for Yuhang Yang
% Compile: latexmk -xelatex -gg -interaction=nonstopmode -halt-on-error yuhang_yang_cv.tex
% ====================================================================
\documentclass[11pt]{article}

\usepackage{cv}

\cvname{阳宇航}
\cvheadline{南京大学数学学院 · 信息与计算科学专业本科生}
\cvphone{(+86) 193 5846 6910}
\cvlocation{Nanjing, China}
\cvemail{yuhangyang@smail.nju.edu.cn}
\cvwebsite[yuhangyang.site]{https://yuhangyang.site}

\begin{document}

\makecvheader

\cvsection{研究兴趣}

当前重点探索计算数学与人工智能的交叉方向，关注数值优化与科学计算方法，以及深度学习和大语言模型的数学基础、训练与推理机制。

\cvsection{教育背景}

\begin{cventry}{2024--2028（预计）}{信息与计算科学专业（本科）}{南京大学数学学院}{中国·南京}
学位学分绩：\textbf{4.45/5.0}；排名：\textbf{1/15}。
\begin{cvitems}
  \item 核心课程：数学分析（85）、高等代数（92）、解析几何（99）、数值分析（91）、离散数学（95）。
  \item 英语能力：大学英语四级（CET-4）627 分；大学英语六级（CET-6）508 分。
\end{cvitems}
\end{cventry}

\cvsection{技术技能}

\cvitem{编程语言}{Python（熟悉）、C++（基础）}
\cvitem{深度学习框架}{PyTorch（基础）}
\cvitem{工具}{\LaTeX（熟悉）、Git（基础）}

\cvsection{竞赛经历}

\cvitem{2025}{\textbf{二等奖}——全国大学生数学竞赛（江苏赛区）}
\cvitem{2025}{\textbf{三等奖}——全国大学生数学建模竞赛（江苏赛区）}

\cvsection{奖项与荣誉}

\cvitem{2025}{\textbf{南京大学人民奖学金二等奖}（专业前 20\%）}

\end{document}
